\subsection{二つの長さ2基点を持つCY射影の全H3排除}
\label{subsec:cubic-CY-finite22-cover}
\begin{theorem}\label{thm:cubic-CY-finite22-empty}
零交換子CY包絡を持つH3の対称射影の基点型は22にならない。
従って残る必要射影型は1111と211である。
\end{theorem}
\begin{proof}
固定出典\cite{AS4CYpinned}の二つの長さ2基点の必要被覆
(TT1)--(TT14)を用いる。AS正則性の十分性や全自己同型群は用いない。
対称射影は $U=(b^2,d^2,ac,ad,bc,bd)$、印付き関係は
$Q=-z_2z_5+z_3z_4-z_0z_1+z_5^2$ とする。
全安定化群は基点環 $k[t]/t^2\times k[s]/s^2$ の全切断への作用から得られ、
二成分の交換と
\[
 N(p,q)=\begin{pmatrix}1&0&0&0\\p&1&0&0\\0&0&1&0\\0&0&q&1\end{pmatrix},\qquad
 T(\alpha,v)=\operatorname{diag}(1,\alpha,v,v/\alpha)
\]
で生成される。混合CY核は全21係数であり、元関係は
$R=U+KB^{\mathsf T}A^{-1}$、残余条件は
$C+gJ-B^{\mathsf T}A^{-1}B=0$ である。
これらの行列は公開検証器と最小固定入力に保持する。

$\kappa=x_6\ne0$ なら $N$ の作用
$x_7'=x_7+q\kappa$、$x_{12}'=x_{12}+p\kappa$ により両係数を消せる。
$x_0$ 又は $x_5$ 非零の枝は $\mathcal A$ に尽くされる。
両方零なら、$x_{10}$ 又は $x_{11}$ 非零の枝の一つは空idealであり、
残りは $\kappa=\pm i/8$ の二族で四次商次元38。
両係数零なら $\kappa=\pm1/4$ の二代表で四次商次元40となり、全て排除される。
ここで使う消去の分母は有理定数と明記した非零係数の冪だけである。

$\kappa=0$、$(x_7,x_{12})\ne(0,0)$ なら交換とトーラスで $x_7=1$、
$N$ の定数pivot 2と1で $x_8=x_{16}=0$ とする。
$x_0\ne0$ の枝は $\mathcal B$ 又は四次商次元39の一族。
$x_0=0,x_5\ne0$ は実平方 $d^2$ を含む。
$x_0=x_5=0$ では残余式 $x_9^2=1$ と
$x_2(x_9+1)=x_{15}(x_9-1)=0$ を用い、
平方関係の枝を除くと、四次商次元42の一族又は三次商次元21の一族になる。
従ってこの場合も $\mathcal B$ だけが残る。

$\kappa=x_7=x_{12}=0$ では $x_{10}x_{11}=0$。
どちらか非零なら交換後 $x_{11}=1$ とでき、定数残余式 $1/8=0$ に矛盾する。
従って両方零。$x_0$ 又は $x_5$ 非零なら交換後 $x_0=1$ とし、
$x_5=1/16-r^2$ を得る。
$x_5=0$ は実平方 $b^2$、$r\ne0$ は $\mathcal C_r$、$r=0$ は $\mathcal D$ に移る。
最後に $x_0=x_5=0$ とし、
$\delta=x_{17}-x_9$、$\epsilon=x_{14}-x_{18}$ と置くと
\[
 \delta\epsilon=1/4-4r^2,\qquad
 rD_{24}-x_{14}D_{22}/2=2r^2\delta.
\]
右辺の判別値が非零なら $r=0$ が従い、
$\delta=t\ne0$ として $\mathcal E_t$ 又は $t=-2$ の $\mathcal F_h$ になる。
平行移動の係数 $t+2$ が消える後者では、自由係数 $h$ を残す。
判別値零の枝は $r=\pm1/4$ であり、平方を避けた全残余が三次商次元24の二族になる。

以上の九つの非正則な代表・族の元商次元は、
当repoで元六関係から独立に再計算した。
族の分母は明記した非零自由係数の冪だけである。
その環上の全語関係行列のrankは特殊化で増えないので、
一般商次元の下界が全許容閉値へ及ぶ。
従って例外値を落とさず全て排除できる。

残る六表示は $\mathcal A,\mathcal B,\mathcal D,\mathcal C_r,\mathcal E_t,\mathcal F_h$。
元16語上の六関係行列が既存の各表示と文字どおり一致することを確認した。
$\mathcal A$ の両符号も含む。
定理\ref{thm:tt22-C-all-excluded}、定理\ref{thm:tt22-EF-all-excluded}、
命題\ref{prop:tt22-fixed-all-excluded}から、全てH3包絡にならない。
\end{proof}

独立検証器は\\
\texttt{verify\_cubic\_CY\_finite22\_cover.py}。
全21混合核、印付き群の共変性、全九非正則商と六表示の移送を確認する。
全分岐正規化には固定出典の通常証明と、隔離した上流全control再実行を用いる。
必要schemeの被約性、族の相対被覆、全CY分類や新規性を同時に主張しない。
